\lecture{11}
\title{PAC and VC Dimension}

\begin{document}

\subsection{Probably Approximately Correct (PAC)}

The result from the end of
\href{/6.790/lectures/10/#estimation-error-bounding}{Lecture 10}
inspires the following definition.

\textbf{Definition.}
For $\epsilon > 0$ and $\delta \in (0, 1)$,
a learning algorithm that takes training data $\mathcal{D}$
and outputs a hypothesis $\hat{h}$ is
$(\epsilon, \delta)$-\textbf{PAC} if:
\[ \mathbb{P}\left(
    \mathcal{R}(\hat{h}) - \mathcal{R}(h^*) \leq \epsilon
\right) \geq 1 - \delta, \]
where, as before, we denote
$h^* := \argmin_{h \in \mathcal{H}} \mathcal{R}(h)$.
\begin{itemize}
    \item Think of $\epsilon$ as an accuracy parameter
    that represents the difference between
    the learned and optimal hypotheses $\hat{h}$ and $h^*$.
    \item Think of $\delta$ as a confidence parameter
    that represents the chance that our data $\mathcal{D}$
    misrepresents the actual data distribution.
\end{itemize}

\textbf{Example.}
The result from
\href{/6.790/lectures/10/#estimation-error-bounding}{Lecture 10}
says that, if the risk function $\mathcal{R}(h)$
has a loss function satisfying $\mathcal{L}(x, y) \in [0, 1]$,
then the ERM learning algorithm
is $\left(2 
\sqrt{\frac{\log |\mathcal{H}| + \log\left(2 / \delta\right)}{2|\mathcal{D}|}},
\delta \right)$-PAC.

\subsection{PAC-Learnable}

\textbf{Definition.}
A hypothesis class $\mathcal{H}$ is \textbf{PAC-learnable}
if there exists a learning algorithm
that takes training data $\mathcal{D}$
and outputs a hypothesis $\hat{h}$
such that the following is true
for any $\epsilon > 0$ and $\delta \in (0, 1)$.
\begin{quote}
    There exists a constant $m_{\mathcal{H}}(\epsilon, \delta)$
    such that, no matter what the data distribution,
    $\hat{h}$ is $(\epsilon, \delta)$-PAC
    whenever $|\mathcal{D}| > m_{\mathcal{H}}(\epsilon, \delta)$.
\end{quote}

\textbf{Example.}
The result from
\href{/6.790/lectures/10/#estimation-error-bounding}{Lecture 10}
says that any \emph{finite} hypothesis class $\mathcal{H}$
is PAC-learnable (assuming a bounded loss function).
In particular, in this case we have
$m_{\mathcal{H}} = O\left(
    \frac{\log|\mathcal{H}| + \log(1/\delta)}{\epsilon^2}
\right)$.

The assumption of ``bounded loss'' in the above example
actually turns out to be pretty much required.

\begin{quote}
    \textbf{Fact.}
    If the loss function is unbounded,
    then it turns out almost \emph{nothing} is PAC-learnable.

    \emph{Justification:}
    Let's consider the simplest nontrivial case: $\mathcal{H} = \{h_1, h_2\}$.
    Then consider the following distribution on $\mathcal{D}$.
    \begin{enumerate}
        \item[(a)] With very small probability $p$,
        we sample $(x, y)$ for which
        $\mathcal{L}(h_1(x), y) = M$
        and $\mathcal{L}(h_2(x), y) = 1$,
        where $M \gg 1$.
        \item[(b)] With very large probability $1 - p$,
        we sample $(x, y)$ for which
        $\mathcal{L}(h_1(x), y) = \mathcal{L}(h_2(x), y)$.
    \end{enumerate}
    Then no matter how large an $m_{\mathcal{H}}(\epsilon, \delta)$ we pick,
    we can set $p$ small enough so that
    $\mathcal{D}$ almost always hits case (b).
    Then it is impossible for the learning algorithm
    to learn that $h_2$ is the better hypothesis.

    Thematically, the reason why unbounded loss
    leads to PAC-unlearnability is that
    \emph{unbounded} nonnegative random variables
    can have means that are
    impossible to bound
    given only finitely many samples.
\end{quote}

The conclusion, then, is that
we might as well assume the loss function is bounded.
We should also direct our attention to
\emph{infinite} hypothesis classes $\mathcal{H}$,
as all the finite ones are automatically PAC-learnable
under this assumption.

\subsection{Growth Functions: Motivation}

From now on,
we'll specifically consider the case when
our hypotheses $h \in \mathcal{H}$
are binary classifiers
$h \colon \mathcal{X} \to \{+1, -1\}$.
Even still, proving or disproving PAC-learnability
for infinite hypothesis classes seems hard.

\begin{remark}
    This subsection and the next
    will take a step back from PAC-learnability for a while.
    Just have trust that
    all of this setup will eventually prove itself
    incredibly useful for deciding
    whether an infinite hypothesis class
    is PAC-learnable or not.
\end{remark}

Suppose $m$ is the size of our training data
$\mathcal{D} = \{(x_i, y_i)\}_{i = 1}^m$.
Then there are only so many ways
a hypothesis $h \in \mathcal{H}$ can perform
in the training process.
The training data acts like
an $m$-question true-or-false exam,
so there are at most $2^m$ ways
a hypothesis $h \in \mathcal{H}$ could perform.
Our learning algorithm should
pick the hypothesis $\hat{h} \in \mathcal{H}$
that scores the highest.

But if two hypotheses $\{h_1, h_2\}$ score the highest\dots
well, then we haven't learned enough.
It's possible that our learning algorithm outputs $\hat{h} = h_1$,
whereas the true optimal hypothesis is $h^* = h_2$.
And maybe on the real test data,
the output $\hat{h}$ performs horribly
in comparison to $h^*$.
We'd have no way of knowing,
because $\hat{h}$ and $h^*$ performed the same
in the training exam.

\subsection{Growth Functions: Definition}

To formalize the idea of multiple hypotheses
ending up indistinguishable in training,
it helps to \emph{count}
the number of possible ways
a hypothesis $h \in \mathcal{H}$
can perform on our $m$-question true-or-false exam.

\textbf{Definition.}
Given a training set $X = (x^{(1)}, \dots, x^{(m)}) \in \mathcal{X}^m$,
let $\mathcal{H}_X$ denote the set of all possible
labelings of points in $X$ by some hypothesis in $\mathcal{H}$; that is,
\[ \mathcal{H}_X := \left\{
    \left(h(x^{(1)}), \dots, h(x^{(m)})\right)
    \mid h \in \mathcal{H}
\right\}. \]

\textbf{Definition.}
Define $\Gamma_{\mathcal{H}}(m)
:= \displaystyle \max_{X \in \mathcal{X}^m} |\mathcal{H}_X|$
to be the \textbf{growth function} of $\mathcal{H}$.

\textbf{Example.}
In the worst case, we have
$\Gamma_{\mathcal{H}}(m) = 2^m$ for all $m$.
This represents the scenario in which
multiple hypotheses are the highest-scoring,
making it impossible to pick out a PAC hypothesis.

\begin{remark}
    By definition, $\Gamma_{\mathcal{H}}(m) > 2^m$ is impossible.
\end{remark}

\begin{problem}
    Determine the growth function of
    $\mathcal{H} := \{x \mapsto \mathrm{sgn}(x - \theta)
    \mid \theta \in \mathbb{R}\}$.
\end{problem}

\begin{solution}
    Suppose the training set is $x^{(1)} < \dots < x^{(m)}$,
    in that order.  Then there are only
    $\Gamma_{\mathcal{H}}(m) = \boxed{m + 1}$ ways
    a hypothesis in $\mathcal{H}$ could label these points;
    one such way is pictured below.
    \begin{center}
        \frame{\includegraphics[width=0.6\textwidth]{lecture-11/threshold-growth}}
    \end{center}
\end{solution}

\begin{problem}
    Let $\mathcal{H}$ be the class of all
    two-dimensional linear classifiers.
    Compute $\Gamma_{\mathcal{H}}(3)$ and
    $\Gamma_{\mathcal{H}}(4)$.
\end{problem}

\begin{solution}
    The picture below demonstrates that
    $\Gamma_{\mathcal{H}}(3) = \boxed{8}$
    and is achieved by
    any three points in general position.
    \begin{center}
        \frame{\includegraphics[width=0.3\textwidth]{lecture-11/linear-3}}
    \end{center}
    Interestingly, $\Gamma_{\mathcal{H}}(4)$ does not equal $2^4 = 16$.
    If the four points of our data set are taken to be
    the vertices of a square, then
    not all $2^4$ labelings of those vertices are
    achievable by a linear classifier.
    \begin{center}
        \frame{\includegraphics[width=0.1\textwidth]{lecture-11/linear-4}}
    \end{center}
    Infamously, as shown above,
    the XOR cannot be modeled by a linear classifier,
    and $\Gamma_{\mathcal{H}}(4) = \boxed{14}$ instead of $16$.
    (To prove this, we would have to check that
    no other arrangement of four points
    gives more than $14$ achievable labelings.)
\end{solution}

\subsection{VC Dimension}

Shockingly, as the following theorem states,
there are only two types of growth functions.

\begin{theorem}{VC Dichotomy}
    For any hypothesis class $\mathcal{H}$, either:
    \begin{itemize}
        \item We have $\Gamma_{\mathcal{H}}(N) = 2^N$
        for all $N$.
        \item We have $\Gamma_{\mathcal{H}}(N) \leq C \cdot N^d$
        for all $N$, for some constants $C > 0$ and $d \in \mathbb{N}$.
    \end{itemize}
\end{theorem}

\begin{remark}
    Note that this theorem is fundamentally
    a set-theoretic result!
    Nothing about its statement
    requires any mention of a probability distribution
    or machine learning at all.
\end{remark}

This inspires the following definitions.

\textbf{Definition.}
A training set $X \in \mathcal{X}^m$ is \textbf{shattered}
by a hypothesis class $\mathcal{H}$ if
$|\mathcal{H}_X| = 2^m$.
The \textbf{VC-dimension} $\mathrm{VC}(\mathcal{H})$
of a hypothesis class $\mathcal{H}$
is the size of the largest subset of $\mathcal{X}$
that is shattered by $\mathcal{H}$.

\textbf{Example.} If $\mathcal{H}$ is such that
$\Gamma_{\mathcal{H}}(N) = 2^N$ for all $N$, then
its VC-dimension is infinite.  Meanwhile,
\begin{itemize}
    \item The VC-dimension of the hypothesis class
    $\mathcal{H} := \{x \mapsto \mathrm{sgn}(x - \theta)
    \mid \theta \in \mathbb{R}\}$ from earlier is $1$.
    \item The VC-dimension of the class of
    all two-dimensional linear classifiers is $3$.
    \item More generally, you can check that
    the VC-dimension of the class of all
    $d$-dimensional linear classifiers is $d + 1$.
\end{itemize}

\begin{remark}
    You might notice that
    the VC-dimension of each of these examples
    happens to equal the number of parameters
    associated with its hypothesis class.
    Please note that this does \textbf{not} hold in general.
\end{remark}

\subsection{Results: Identifying PAC-Learnability}

Now we can finally get back to PAC-learnability!
Let's state (without proof) some results using our newfound definitions.

\begin{theorem}{ERM VC Risk Bound}
    Under the ERM learning algorithm,
    the following holds with probability at least $1 - \delta$.
    \[
    \mathcal{R}(\hat{h}) \leq
    \mathcal{R}(h^*) + 8 \cdot \sqrt{
        \frac{\log\left(\Gamma_{\mathcal{H}}(2|\mathcal{D}|)\right)
        + \log(2/\delta)}{2|\mathcal{D}|}
    }.
    \]
\end{theorem}

It also turns out that the exponent $d$
in the bound $\Gamma_{\mathcal{H}}(N) \leq C \cdot N^d$
relates to the VC-dimension!

\begin{theorem}{Sauer-Shelah-Perles}
    If a hypothesis class $\mathcal{H}$
    has finite VC-dimension $d := \mathrm{VC}(\mathcal{H})$, then
    for all $N$ we have:
    \[ \Gamma_{\mathcal{H}}(N) \leq \sum_{i = 0}^d \binom{N}{i}. \]
    In particular, this implies that
    $\Gamma_{\mathcal{H}}(N) = O(N^d)$,
    so the degree of polynomial growth of $\Gamma_{\mathcal{H}}(N)$
    is bounded by $\mathrm{VC}(\mathcal{H})$.
\end{theorem}

\begin{remark}
    Again, the above theorem is purely set-theoretic.
    It upper-bounds the size of a collection of subsets
    $S_i \subseteq \{1, \dots, N\}$ for which 
    no $d + 1$ elements of $\{1, \dots, N\}$ are shattered.
    (The 6.790 notes offer a proof, but I believe it is a fakesolve.)
\end{remark}

Putting the previous two theorems together
yields the \textbf{Fundamental Theorem of Statistical Learning}.

\begin{theorem}{Fundamental Theorem}
    A hypothesis class $\mathcal{H}$ is PAC-learnable
    if and only if it has finite VC-dimension.

    If its VC-dimension $d := \mathrm{VC}(\mathcal{H})$ is finite,
    then, with probability at least $1 - \delta$, we furthermore have:
    \[
    \mathcal{R}(\hat{h}) \leq \mathcal{R}(h^*) + O\left(
        \sqrt{\frac{d}{|\mathcal{D}|} \log \frac{|\mathcal{D}|}{d}
        + \frac{1}{|\mathcal{D}|} \log \frac{2}{\delta}}
    \right).
    \]
    Expectedly, the RHS tends to zero as $|\mathcal{D}|$ grows large,
    confirming PAC-learnability.
\end{theorem}

Yippee!  And so all those hypothesis classes
we discussed earlier are now proven PAC-learnable.

\end{document}